% ============================================================================
% Section 6.5 — Cross-Sections of Three-Dimensional Figures
% CCSS 7.G.A.3
% ============================================================================
\section{Cross-Sections of 3D Figures}

\begin{stepGoal}
\begin{itemize}
    \item Describe the 2D shape formed when a plane slices a 3D figure
    \item Identify cross-sections of prisms and pyramids
\end{itemize}
\end{stepGoal}

\begin{stepVocabBox}
    \vocabItem{Cross-Section}{The flat (2D) shape you see when you slice through a 3D figure with a plane.}
    \vocabItem{Plane}{A flat surface that extends in all directions, like a sheet of glass cutting through the solid.}
\end{stepVocabBox}

\begin{stepsCard}[How to Find the Cross-Section of a 3D Figure]
    \stepItem{Identify the \textbf{3D figure} (prism, pyramid, cylinder, cone, sphere).}
    \stepItem{Determine the \textbf{direction of the cut}: parallel to the base, perpendicular, or diagonal.}
    \stepItem{Name the \textbf{2D shape} that results: rectangle, triangle, trapezoid, circle, etc.}
\end{stepsCard}

% --- Worked Example 1 ---
\begin{stepExample}{You slice a rectangular prism (box) with a cut parallel to its base. What shape do you get?}
    \stepShow{1} The figure is a \textbf{rectangular prism}.

    \stepShow{2} The cut is \textbf{parallel to the base} (horizontal).

    \stepShow{3} A horizontal slice of a rectangular prism has the same shape as the base.
    \stepResult{The cross-section is a rectangle.}
\end{stepExample}

% --- Worked Example 2 ---
\begin{stepExample}{A rectangular pyramid is sliced horizontally halfway up. What shape appears?}
    \stepShow{1} The figure is a \textbf{rectangular pyramid}.

    \stepShow{2} The cut is \textbf{parallel to the base}, halfway between the base and the apex.

    \stepShow{3} A horizontal cut through a rectangular pyramid gives a smaller version of the base shape.
    \stepResult{A smaller rectangle.}
\end{stepExample}

\mascotStepTip{Every horizontal slice of a prism looks exactly like its base. Pyramids narrow, so the slice is a smaller copy!}

% ============================================================================
% PRACTICE
% ============================================================================
\newpage
\begin{stepPractice}{Cross-Sections Practice}
    \resetProblems

    \stepPracticeHeader[step1Color]{Name the Figure}
    \begin{multicols}{2}
    \prob A soup can is a \answerBlank[2.5cm].
    \answer{cylinder}
    \prob A box of cereal is a \answerBlank[2.5cm].
    \answer{rectangular prism}
    \end{multicols}

    \stepPracticeHeader[step2Color]{Direction of the Cut}
    \prob You slice a cube diagonally from one edge to the opposite edge. What kind of cut is this? \answerBlank[2.5cm]
    \answer{Diagonal}

    \stepPracticeHeader[step3Color]{Name the Cross-Section}
    \prob Horizontal cut through a cone, parallel to its base. Shape? \answerBlank[2.5cm]
    \answer{Circle}

    \prob Vertical cut through a rectangular prism, parallel to a face. Shape? \answerBlank[2.5cm]
    \answer{Rectangle}

    \wordProblem{Is every cross-section of a sphere a circle? Explain.}{~}
    \answer{Yes. No matter how you slice a sphere, the cross-section is always a circle. The size depends on how close the cut is to the center.}
\end{stepPractice}
